\documentclass[11pt]{article}

% --- Packages ---
\usepackage[margin=1in]{geometry}
\usepackage{amsmath,amssymb,amsfonts}
\usepackage{graphicx}
\usepackage{booktabs}
\usepackage{natbib}
\usepackage{hyperref}
\usepackage{xcolor}
\usepackage{enumitem}
\usepackage{bm}
\usepackage{float}
% subcaption not needed: 3-panel compositing done in Python

% --- Custom commands ---
\newcommand{\vect}[1]{\bm{#1}}
\newcommand{\mat}[1]{\mathbf{#1}}
\DeclareMathOperator{\softplus}{softplus}
\DeclareMathOperator{\sigmoid}{sigmoid}
\DeclareMathOperator{\logit}{logit}

\title{Six-hourly Probabilistic Precipitation Forecasts Created from Hourly Probabilities}

\author{Thomas M.\ Hamill\thanks{The Weather Company, Atlanta,
  Georgia; e-mail: \texttt{tom.hamill@weather.com}}}

\date{\today}

\begin{document}

\maketitle

\begin{abstract}
A lightweight multi-layer perceptron (MLP) is described that produces
six-hourly, spatially detailed probabilistic quantitative precipitation forecasts (PQPF) from
a sequence of hourly spatially detailed probabilistic forecasts.  The hourly forecasts were
generated by an Attention Residual U-Net that outputs the parameters of a
probability distribution including a probability of no precipitation and a weighted 
mixture of two Gamma distributions at each grid point for each hourly
forecast lead time.  The MLP takes as input the six distributional parameters
from each of six consecutive hours and maps them to a corresponding
six-parameter Gamma mixture describing the six-hourly accumulated precipitation.
Training uses importance-sampled pairs of hourly probabilistic forecasts 
generated from the Attention Residual U-Net and
matched Multi-Radar Multi-Sensor (MRMS) six-hourly precipitation
accumulations, subject to a quality threshold.  The resulting six-hourly probabilistic 
forecasts are well calibrated and skillful, as demonstrated by reliability and performance 
diagrams evaluated over an out-of-sample period spanning day 10 through the
end of every month of 2025.
Forecasts show realistic modulation of precipitation related to terrain variations.
\end{abstract}

% =====================================================================
\section{Introduction}
\label{sec:intro}
% =====================================================================

There is now a substantial body of literature on methods for the statistical postprocessing of
ensemble forecasts of precipitation for accumulation periods of 6, 12, or 24 hours.  See, for example,
\citet{hamill2023nbm} and \citet{jesberg2023}, who describe methods and results for improving
National Blend of Models six-hourly probabilistic precipitation forecasts using long time series
of reforecasts and precipitation reanalyses.  Other past examples include
\citet{sloughter2007}, \citet{voisin2010}, \citet{bentzien2012}, \citet{bouballegue2013},
\citet{scheuerer2014}, \citet{messner2014}, \citet{scheuerer2015}, \citet{baran2016},
\citet{zhang2017}, \citet{hamill2017nbm}, \citet{hamill2018}, \citet{gascon2019},
\citet{ghazvinian2020}, and \citet{ghazvinian2022}.

In a previous article, a method was described for how one may
generate hourly high-resolution probabilistic precipitation forecast guidance from
a deterministic weather forecast model simulation through an Attention
Residual U-Net deep learning algorithm \citep{hamill2026a}.
These hourly probabilities may be useful for providing context on precipitation event
starting and end times and intensity.  Still, in many situations, a user may want to also have internally
consistent probabilities over longer periods such as a 6-h window (say, the morning
or the afternoon) or a 12-h window (day, night).

When probability forecasts are issued
independently for successive hourly periods, combining them into a multi-hour
accumulation forecast is complicated by the serial correlation of precipitation
events.  Consecutive hours are often not meteorologically independent, and treating
them as such may produce poorly calibrated accumulation probabilities.
\citet{hughes1979} addressed this type of problem, combining
routine 12-h probability-of-precipitation forecasts from National Weather Service
offices into longer-period estimates.  Hughes and Sangster applied empirical joint-probability
adjustments for forecast-period dependence and find that warm-season events
behave more nearly independently than cold-season ones.
\citet{wilks1990} examined the combination of forecast probabilities for
consecutive precipitation periods more generally and shows that the joint
distribution of successive events, not simply the product of marginal
probabilities, must be accounted for to achieve reliable aggregated guidance.
\citet{wilks1991} further shows how such serial correlation of meteorological
events and forecasts can be represented explicitly within dynamic
decision-analytic models, reinforcing that treating consecutive periods as
independent is a simplification with real decision-relevant consequences.

Papers such as these illustrate the problem of serial correlation but do not
provide an actionable algorithm for generating, say, six-hourly probabilities
from hourly data.  We thus seek an alternative, inspired by recent success
with deep learning models for precipitation statistical postprocessing
\citep[e.g.,][]{hu2023}.
The approach to be described trains a model to aggregate the distributional information
carried by six consecutive hourly forecasts into a single six-hourly predictive
distribution.  In principle, the same approach could be used for other accumulation
periods such as 3 h, 12 h, or 24 h, though they will not be demonstrated here.

This document describes this trained model, a lightweight multi-layer perceptron
(MLP; \citealt{rumelhart1986}) that maps the 38-dimensional feature vector formed
by six consecutive hourly gamma-mixture parameter sets plus two cyclic
day-of-year features to the six parameters of a
six-hourly zero-inflated two-component Gamma mixture, i.e. a probability mass at zero
and a weighted Gamma mixture model for positive precipitation amount.  The MLP is trained separately
for each lead time using importance-sampled pairs of hourly probabilistic forecasts
and matched MRMS six-hourly precipitation sums.  By operating directly on the
distributional parameters of the hourly forecasts rather than on raw precipitation
fields, the approach avoids the need to store or process full-resolution spatial
fields at training time, and the resulting six-hourly probabilistic forecasts
inherit the calibration properties of the underlying hourly model.

The remainder of this document is organized as follows.
Section~\ref{sec:hourly} describes the input data, including the probabilistic
forecast feature data, the MRMS target precipitation, and the sampling strategy
used to construct the training set.
Section~\ref{sec:mlp} describes the MLP architecture, loss function, and training.
Section~\ref{sec:results} presents reliability diagrams for 12-h and 48-h lead-time
forecasts.
Section~\ref{sec:summary} provides a summary and possible next steps.

% =====================================================================
\section{Input data}
\label{sec:hourly}
% =====================================================================

\subsection{Probabilistic forecast feature data}
\label{sec:input_feat}

The hourly probabilistic forecasts described by \citet{hamill2026a}, used as input
to the MLP, are produced by an Attention Residual U-Net (Attention ResUNet) applied
to The Weather Company's GRAF (Global high-Resolution Atmospheric Forecasting) model
output.  GRAF is a convection-permitting model based on the NCAR Model for Predictions
Across Scales \citep[MPAS;][]{skamarock2012}, with mesh refinement to approximately 4.25\,km
grid spacing over the contiguous United States (CONUS) and Europe.

At each grid point and for each hourly lead time, the Attention ResUNet predicts
the six parameters of a zero-inflated two-component Gamma mixture:
\begin{equation}
  \theta_h = (p_0,\; w,\; \alpha_1,\; \theta_1,\; \alpha_2,\; \theta_2)_h,
  \label{eq:hourly_params}
\end{equation}
where $p_0$ is the probability of zero hourly precipitation, $w$ is the
conditional mixing weight for the lighter component, and
$(\alpha_j, \theta_j)$ are the shape and scale parameters of the $j$-th
Gamma component.  Component 1 is constrained to represent lighter
precipitation than component 2.

The probability of exceeding any precipitation amount threshold $\tau > 0$ is obtained
analytically from the mixture cumulative distribution function.  The
parameters $\theta_h$ are saved for each grid point, initial condition date and lead time,
and these files serve as the primary data source for both the MLP training
and the MLP inference pipelines described below.

The six distributional parameter fields from each of the six hourly forecasts
are stacked, together with two cyclic day-of-year features encoding the
sample's annual season, to form a 38-dimensional feature vector per
training sample:
\begin{equation}
  \vect{x} = \bigl[
    p_{0,1},\ldots,p_{0,6},\;
    w_1,\ldots,w_6,\;
    \alpha_{1,1},\ldots,\alpha_{1,6},\;
    \theta_{1,1},\ldots,\theta_{1,6},\;
    \alpha_{2,1},\ldots,\alpha_{2,6},\;
    \theta_{2,1},\ldots,\theta_{2,6},\;
    \cos\!\left(\tfrac{2\pi d}{365}\right),\;
    \sin\!\left(\tfrac{2\pi d}{365}\right)
  \bigr] \in \mathbb{R}^{38},
  \label{eq:feature_vec}
\end{equation}
where the second index in the vector elements above corresponds to the
six consecutive hourly lead times, ordered from earliest to latest, and
$d$ is the initial condition's day of year (1--365).  The cyclic day-of-year
pair gives the MLP an explicit seasonal signal, which it otherwise has no
way to infer from the hourly distributional parameters alone.

\subsection{Target MRMS data}
\label{sec:mrms}

The training target at each sampled grid point is the six-hourly Multi-Radar
Multi-Sensor \citep[MRMS;][]{zhang2016} precipitation accumulation, computed 
as the sum of the six matched hourly MRMS analyses:
\begin{equation}
  y = \sum_{k=1}^{6} r_k, \quad y \geq 0 \text{ mm},
  \label{eq:target}
\end{equation}
where $r_k$ is the MRMS hourly precipitation for the $k$-th lead time.

Grid points are included in the training set only if their mean MRMS
data-quality index over the six verification hours exceeds a threshold
$Q = 0.6$.  This is more restrictive than the quality threshold used for
training the hourly AttnResUNet ($Q > 0.01$), because the six-hourly
accumulation compounds any retrieval uncertainty across six hours; a higher
quality threshold reduces the risk of training on accumulations that are
substantially influenced by radar-beam blockage or other retrieval problems.

\subsection{Sampling strategy}
\label{sec:sampling}

Training data are assembled from previously generated
hourly probability distribution parameters.  Forecasts are
sampled from a selected subset of model grid points, described below, which represent a small
fraction of the overall number of grid points in the domain.  Training and verification use the CONUS domain only (the Europe domain
mentioned in Sec.~\ref{sec:input_feat} is not used here).  Two thousand
sample points are drawn from a gridded forecast with a domain size of
(1524, 1308), which results in approximately
1/1000th of the domain's grid points being selected for training samples.  Samples are drawn
from the first seven calendar days of every month of 2025, with draws from forecasts
initialized at 00 UTC and 12 UTC.  Spanning all twelve months lets the MLP learn genuine
seasonal variation via the cyclic day-of-year features of
Eq.~\eqref{eq:feature_vec}, rather than only the handful of months it would
otherwise be trained on.  Days 8 and 9 of each month are withheld entirely as a
gap, and days 10 through the end of the month are reserved for out-of-sample
testing (Sec.~\ref{sec:results}); the gap prevents a persistent synoptic
system from straddling the training/testing boundary and inflating apparent
skill.  Fitted parameters of the zero-inflated two-component
Gamma mixture are extracted for each of six individual hours preceding forecast lead time.
For example hourly forecast parameters for hours 19, 20, 21, 22, 23, and 24 are extracted
when training the 24-h accumulated precipitation forecast.

Because most grid points are dry at any given time, uniform random sampling
would yield a severely imbalanced training set.  Instead, each grid point is
assigned a sampling weight proportional to its likelihood of being wet.
Specifically, the minimum fraction-zero probability across the six hourly
lead times is used to estimate the probability of positive precipitation
at that point during the window:
\begin{equation}
  p_{\text{samp}} = a + (1-a)(1 - \min_k p_{0,k}),
  \label{eq:importance}
\end{equation}
where $a = 0.01$ is a floor that ensures all quality-passing grid points
retain a small probability of being selected.  The minimum (rather than,
e.g., the product of the six $p_{0,k}$ values, which would estimate the
probability of at least one wet hour under an independence assumption) is
used because it identifies the single most-likely-wet hour in the window
directly from the forecast itself, without invoking the independence
assumption that Sec.~\ref{sec:results} explicitly contrasts the MLP
against.  Two thousand samples
per initial condition are drawn without replacement according to the normalized
$p_{\text{samp}}$ weights, subject to the MRMS quality criterion being met.
Samples are concatenated over all twelve months, each
sample comprising 36-dimensional hourly feature arrays
(shape $N \times 6$ per variable) plus the two cyclic day-of-year features,
the six-hourly MRMS target precipitation,
the mean MRMS quality, and the latitude and longitude of each sampled point.

% =====================================================================
\section{MLP Architecture, Loss Function, and Training}
\label{sec:mlp}
% =====================================================================

\subsection{Architecture}
\label{sec:architecture}

Six-hourly probabilistic forecasts are produced using a fully connected
multi-layer perceptron.  The network maps the 38-dimensional
input feature vector~\eqref{eq:feature_vec} to six unconstrained real
values, which are subsequently transformed to the constrained parameters
of a zero-inflated two-component Gamma mixture.

The hidden layer widths are $[72, 144, 72, 36, 12]$, giving the full
architecture
\begin{equation}
  38 \to 72 \to 144 \to 72 \to 36 \to 12 \to 6.
  \label{eq:mlp_arch}
\end{equation}
Each hidden layer applies a linear transformation followed by batch
normalization \citep{ioffe2015} and a ReLU activation \citep{nair2010}.  The bottleneck
design, expanding to 144 units before compressing back to 12, encourages
the network to first build a richer joint representation of the six hourly
parameter sets before projecting to the six-hourly distributional summary.
The resulting network has 27{,}582 trainable parameters (linear-layer
weights and biases, plus batch-normalization scale and shift), justifying
the description of the architecture as lightweight: for comparison, the
hourly Attention ResUNet that supplies its input features has approximately
$3.3 \times 10^7$ parameters, roughly 1200 times as many.

Figure~\ref{fig:mlp_arch} provides a schematic of the architecture.

\begin{figure}[H]
\centering
\includegraphics[width=\textwidth]{Figs_6h/mlp_architecture_6h.png}
\caption{Architecture of the MLP neural network for six-hourly probabilistic
  precipitation forecasting.  The network maps 38 input features (six distributional
  parameters from each of six consecutive hourly forecasts, ordered from earliest
  to latest lead time, plus two cyclic day-of-year features) through five hidden
  layers ($72 \to 144 \to 72 \to 36 \to 12$
  units, each followed by batch normalization and ReLU) to six output logits.
  The logits are transformed via sigmoid (for $p_0$ and $w$) and softplus with
  a mean-ordering constraint (for the Gamma shape and scale parameters) to produce
  the six parameters of the six-hourly zero-inflated two-component Gamma mixture.}
\label{fig:mlp_arch}
\end{figure}

\subsection{Output Model Parameters}
\label{sec:output}

The six raw logits $(z_0,\ldots,z_5)$ from the output linear layer are
transformed to the six constrained distributional parameters using the same
activation scheme as the hourly AttnResUNet:
\begin{align}
  p_0 &= \sigmoid(z_0), \\
  w   &= \sigmoid(z_1), \\
  \alpha_1 &= \alpha_{\min} + \softplus(z_2), \\
  \theta_1 &= \theta_{\min} + \softplus(z_3), \\
  \alpha_2 &= \alpha_1 + \softplus(z_4) + \Delta_{\min}, \\
  \theta_2 &= \theta_{\min} + \softplus(z_5),
\end{align}
where $\alpha_{\min} = 0.1$, $\theta_{\min} = 0.01$, and $\Delta_{\min} = 0.5$.
To resolve label switching between the two Gamma components, component~2's
shape parameter is reparameterized as a hard structural offset from
component~1's shape, $\alpha_2 = \alpha_1 + \softplus(z_4) + \Delta_{\min}$,
rather than by sorting the two components after the fact.  This guarantees
$\alpha_2 > \alpha_1$ by construction for every input, so component~2 is
always the larger-shape component and no post-hoc reordering step -- and
therefore no discontinuous, non-differentiable relabeling operation -- is
needed during training.  This mirrors the equivalent hard-ordering constraint
used in the hourly Attention ResUNet's mixture loss function.

\subsection{Loss Function}
\label{sec:loss}

Training minimizes the negative log-likelihood (NLL) of the zero-inflated
two-component Gamma mixture, identical in form to the loss used for the hourly
Attention ResUNet.  For a sample with target six-hourly accumulation $y$:

\paragraph{Case $y = 0$:}
\begin{equation}
  \text{NLL} = -\log(p_0 + \epsilon),
\end{equation}
where $\epsilon = 10^{-7}$ prevents numerical overflow.

\paragraph{Case $y > 0$:}
\begin{equation}
  \text{NLL} = -\log(1 - p_0 + \epsilon)
    - \log\!\bigl[
        w\, f_{\Gamma}(y;\alpha_1,\theta_1)
        + (1-w)\, f_{\Gamma}(y;\alpha_2,\theta_2)
      \bigr],
\end{equation}
evaluated in log-space using the log-sum-exp trick for numerical stability.
Per-sample NLL values (a unitless loss, not a precipitation amount) are
clamped to a maximum value of 100 before averaging across the mini-batch,
preventing single extreme observations from dominating the gradient.

The NLL was preferred over the Continuous Ranked Probability Score (CRPS),
the other common strictly proper scoring rule used in probabilistic
post-processing.  For a zero-inflated two-component Gamma mixture, the CRPS
has no closed-form expression and must be approximated by numerical integration
or Monte Carlo sampling of the mixture CDF at many threshold values---an
expensive inner loop for every sample in every mini-batch.  In contrast, the
Gamma log-PDF has an analytic form (computed via python pytorch.lgamma), so the
NLL can be evaluated exactly in a single vectorized pass.  Minimizing the NLL
is also equivalent to maximum-likelihood estimation of the distributional
parameters, providing a direct link to the parametric model being fitted.

\subsection{Feature Normalization}
\label{sec:normalization}

Each of the 38 input features is independently standardized to zero mean
and unit variance using statistics computed from the training split:
\begin{equation}
  \tilde{x}_k = \frac{x_k - \mu_k}{\sigma_k},
  \label{eq:zscore}
\end{equation}
where $\mu_k$ and $\sigma_k$ are the training-set mean and standard deviation
of feature $k$.  Features with $\sigma_k < 10^{-8}$ are left unchanged
(division by 1 rather than by near-zero).  The normalization statistics
are saved and applied consistently at inference time.

Z-score normalization is preferred over the fixed-maximum normalization used
for the spatial AttnResUNet features because the distributional parameters
($\alpha, \theta$) span a wide range across lead times and seasons, and their
natural scale is not known \textit{a priori}.

\subsection{Climatological Output-Layer Initialization}
\label{sec:initialization}

The bias of the final linear layer is initialized using climatological
statistics derived from the training target values, so that the network's
first forward pass produces predictions near the marginal climatological
distribution rather than the near-uniform predictions that result from default
(near-zero) weight initialization.

The climatological zero fraction $\hat{p}_0$ is estimated directly from the
training targets, and the bias $b_0 = \logit(\hat{p}_0)$ is set accordingly.
For the wet pixels, a single Gamma distribution is fitted by the method of
moments:
\begin{align}
  \hat{\mu}      &= \overline{y_{\text{wet}}}, \\
  \hat{\sigma}^2 &= \text{Var}(y_{\text{wet}}), \\
  \hat{\alpha}   &= \hat{\mu}^2 / \hat{\sigma}^2, \\
  \hat{\theta}   &= \hat{\sigma}^2 / \hat{\mu},
\end{align}
and the biases for both Gamma components are initialized to these values via
the inverse softplus.  The mixture weight bias is set to zero (implying
$w = 0.5$).  This single-Gamma initialization differs from the EM-based
two-component initialization used by the hourly AttnResUNet; in practice it
produced stable convergence from the first epoch for the smaller MLP, though
this was not directly compared against an EM-based two-component
initialization scheme for this model, so the extent to which the simpler
scheme was necessary, as opposed to merely sufficient, is untested.

\subsection{Training Details}
\label{sec:training}

Training is performed with the Adam optimizer \citep{kingma2015} at a
base learning rate of $\eta = 10^{-3}$.  Mini-batches of 1024 samples are
processed.  Within the day 1--7 training window of each month, training
and validation sets are split by calendar day rather than by a random
shuffle of pooled samples: days 1--4 and 6--7 are used for training, and
day 5 is held out for validation, with both the 00Z and 12Z initializations
of a given day assigned to the same split.  This day-block scheme yields a
more independent validation set than a random 80/20 split, since nearby
samples in time are strongly correlated synoptically, at the cost of a
small validation set (one day per month, 24 initializations total across
the full year) whose NLL estimate is correspondingly noisier than a larger
random split would give; this trade-off was not evaluated quantitatively.  A
\texttt{ReduceLROnPlateau} scheduler (reduction factor 0.5, patience
5 epochs) adapts the learning rate based on validation NLL, and training
halts after 5 consecutive epochs without improvement (early stopping),
up to a maximum of 75 epochs.  Gradients are clipped to a maximum
$\ell_2$ norm of 5.0.

% =====================================================================
\section{Results}
\label{sec:results}
% =====================================================================

Reliability of the six-hourly MLP forecasts is evaluated, populated with forecast
and MRMS observation data for each synoptic cycle initial condition from day 10
through the end of every month of 2025 -- an out-of-sample test period disjoint
from the day 1-7 training window used in Sec.~\ref{sec:sampling}.
The probability of the six-hourly accumulation exceeding a threshold
$\tau$ is then computed analytically from the mixture survival function using
the regularized incomplete gamma function.  Six matched hourly MRMS files
are summed to form the verification six-hourly accumulation, and only grid
points with mean MRMS quality exceeding 0.6 are included in verification,
the same threshold used for the training sample (Sec.~\ref{sec:sampling}).

To provide a reference standard against which the MLP's aggregation skill
can be judged, each reliability diagram also includes a naive
\emph{independence-assumption ensemble control}.  For every verification
case, a 50-member Monte Carlo ensemble is drawn at each grid point from the
zero-inflated two-component Gamma-mixture forecast valid at each of the six
hours spanning the accumulation window.  Corresponding ensemble members are
then summed across the six hours under the (deliberately naive) assumption
that hourly precipitation amounts are statistically independent from hour to
hour.  Six-hourly exceedance probability at each threshold is estimated as
the ensemble relative frequency of the summed accumulation exceeding that
threshold.  Because this control consumes exactly the same hourly marginal
distributions that are fed to the MLP as input features, differences between
the control and the MLP isolate the effect of the temporal correlation
structure the MLP learns from training data, rather than differences in the
underlying hourly forecast quality.

The statistical significance of the MLP-versus-control BSS differences is
assessed following \citet{hamill1999}, who recommend collecting all
spatial samples from a given case into a single per-case score before
testing, since nearby grid points within one synoptic event are strongly
correlated and treating them as independent samples would overstate the
effective sample size, and further find that for probabilistic forecasts a
paired $t$ test or the nonparametric Wilcoxon signed-rank test on those
per-case scores is adequate, without requiring the more elaborate
resampling tests needed for verification statistics computed from
non-probabilistic forecasts.  Accordingly, a single domain-mean Brier Score
is computed for the MLP and for the control at each of the $\sim$500
independent verification cases (case days), and the two paired,
per-case series are compared with both a paired $t$ test and a Wilcoxon
signed-rank test.

Figures~\ref{fig:relia_12h} and~\ref{fig:relia_48h} show reliability diagrams
for three precipitation thresholds at 12-h and 48-h lead times, respectively.
Each figure shows observed frequency of exceedance as a function of forecast
probability, evaluated over the 2025 verification period, for both the MLP
and the independence-assumption control.  Brier Skill Scores (BSS) relative
to the climatological sample base rate are also provided in each panel.
The MLP forecasts are highly reliable, and skill is comparable to or larger than
that reported by \citet{jesberg2023} using NOAA Global Ensemble Reforecast
data, though the verification data was different in that study.

The independence-assumption control is, by contrast, consistently
overconfident.  For both lead times and for all three thresholds, its
reliability curve falls below the diagonal over much of the probability
range, meaning that naively summing independent hourly draws assigns higher
exceedance probabilities than are actually observed.  This is consistent
with real precipitation events being temporally autocorrelated -- wet and
dry hours often cluster together within a synoptic event -- so treating the six
hourly draws as independent manufactures more small, partial accumulations
than the correlated real-world process produces, inflating the estimated
probability of marginally exceeding a threshold.  For example, at the
0.25\,mm threshold and 12-h lead time, cases assigned a forecast probability
of 90\% by the MLP verify 90.9\% of the time, essentially exact, whereas the
control's 90\% bin verifies only 78.6\% of the time; the discrepancy is
larger still at 48-h lead time (94.7\% for the MLP vs.\ 73.1\% for the
control in the 90\% bin).  The MLP's BSS exceeds the control's at every
threshold examined, though the margin is modest at 12-h lead time: 0.536
vs.\ 0.518 (0.25\,mm), 0.459 vs.\ 0.449 (2.5\,mm), and 0.290 vs.\ 0.287
(10\,mm).  The margin is larger at 48-h lead time: 0.423 vs.\ 0.385
(0.25\,mm), 0.330 vs.\ 0.319 (2.5\,mm), and 0.177 vs.\ 0.173 (10\,mm).  A
paired $t$ test and Wilcoxon signed-rank test on the per-case-day Brier
Score series (see above) indicate that the MLP-versus-control difference is
statistically significant ($p < 0.0001$ by both tests) at the 0.25\,mm and
2.5\,mm thresholds at both lead times.  At the 10\,mm threshold, where the
BSS margin is smallest, the two tests disagree: the paired $t$ test does not
reject the null hypothesis of no difference at 12-h lead time ($p = 0.198$)
and is borderline at 48-h lead time ($p = 0.054$), while the Wilcoxon
signed-rank test rejects it at both lead times ($p = 0.0021$ and
$p = 0.0006$, respectively).  This divergence is consistent with the
day-to-day Brier Score differences at the rare 10\,mm threshold being
non-normally distributed (a small number of days with large score
differences and many days with none), which the parametric $t$ test is
less robust to than the nonparametric Wilcoxon test; the 10\,mm result
should therefore be treated as only tentatively significant.  Overall,
this confirms that the MLP is not merely reproducing the hourly marginal
information already present in its inputs, but is learning genuine temporal
correlation structure in six-hourly precipitation accumulation that a naive
independence assumption cannot capture, with the caveat above at the
heaviest, most sparsely sampled threshold -- and, as the performance
diagrams below show, this advantage turns out to be concentrated in
probabilistic calibration rather than in the ability to discriminate
events from non-events.

Results for the 48-h lead time (Fig.~\ref{fig:relia_48h}) show the same
qualitative pattern -- a well-calibrated MLP set against a control that is
overconfident, sometimes markedly so, at high forecast probabilities -- but
with the MLP's BSS advantage over the control larger than at 12-h lead
time, particularly at the lightest threshold.

\begin{figure}[H]
\centering
\includegraphics[width=\textwidth]{Figs_6h/Relia_6h_MLP_MRMS_3panel_2025011000_to_2025123112_12h.png}
\caption{Reliability diagrams for six-hourly MLP probabilistic forecasts (blue
  circles) at 12-h lead time, evaluated over day 10 through the end of every
  month of 2025 (an out-of-sample test period disjoint from the day 1-7
  training window; see Sec.~\ref{sec:sampling}), alongside a naive
  independence-assumption ensemble control
  (red squares; see text).  Panels show the
  reliability of six-hourly exceedance probability for (a) 0.25\,mm,
  (b) 2.5\,mm, and (c) 10.0\,mm.  The diagonal dashed line indicates perfect
  reliability.  The inset histogram in each panel shows the frequency of use for
  each probability bin, for both the MLP (blue) and the control (red).}
\label{fig:relia_12h}
\end{figure}

\begin{figure}[H]
\centering
\includegraphics[width=\textwidth]{Figs_6h/Relia_6h_MLP_MRMS_3panel_2025011000_to_2025123112_48h.png}
\caption{As in Fig.~\ref{fig:relia_12h}, but for 48-h lead time.}
\label{fig:relia_48h}
\end{figure}

Performance diagrams \citep{roebber2009} provide a complementary view of forecast
quality by jointly displaying the probability of detection (POD) and success ratio
(SR, equal to 1 minus the false alarm ratio), with shaded background contours of the
critical success index (CSI) and dashed lines of constant frequency bias.  Curves are
constructed from the same per-bin contingency tables underlying
Figs.~\ref{fig:relia_12h} and~\ref{fig:relia_48h}, by treating each of the 11 saved
forecast-probability bins (0, 10, 20, \ldots, 100\%) as a probability threshold above
which the forecast is counted as ``yes,'' and cumulating hit, false-alarm, and miss
counts from that bin through the highest bin; no additional data collection beyond
what was already gathered for the reliability diagrams is required.
Figures~\ref{fig:perfdiag_12h} and~\ref{fig:perfdiag_48h} show these diagrams for the
12-h and 48-h lead times, respectively, comparing the six-hourly MLP against the
independence-assumption ensemble control at each of the three verification
thresholds.  In contrast to the clear BSS separation reported above, the MLP and
control curves lie close together at both lead times and for all three thresholds,
frequently overlapping or crossing.  Maximum CSI values differ by 0.01 or less between
the two methods at every threshold and lead time examined, and at some
thresholds -- most noticeably the 10\,mm exceedance at 12-h lead time
(Fig.~\ref{fig:perfdiag_12h}c) -- the control's curve sits marginally up and to the
right of the MLP's over part of the probability range.

This is not a contradiction of the BSS results, but a consequence of what each
diagnostic measures.  A performance diagram is built entirely from which
forecast-probability bin a case is placed in, so it reflects discrimination -- the
ability to rank at-risk situations above less-risky ones -- and is insensitive to
whether the numeric probability value attached to a bin is itself trustworthy.
Because the control consumes exactly the same hourly marginal distributions that are
fed to the MLP as input features, the two methods rank cases very similarly, and so
their discrimination, and hence their performance-diagram curves, are nearly
identical.  What differs sharply between the two methods is calibration: as shown
above, the control systematically overstates its own confidence at high-probability
bins, while the MLP's forecast probabilities track the diagonal closely across the
full probability range.  BSS, unlike a performance diagram, penalizes this
miscalibration directly, which is why the MLP's BSS advantage over the control is
real and meaningful even though the two methods' performance diagrams are nearly
indistinguishable.  In short, the MLP's demonstrated advantage over the naive
independence assumption lies in probabilistic calibration rather than in event
discrimination.

\begin{figure}[H]
\centering
\includegraphics[width=\textwidth]{PerfDiagram_6h_MLP_MRMS_3panel_2025011000_to_2025123112_12h.png}
\caption{Performance diagrams (probability of detection vs.\ success ratio, with shaded
  critical success index and dashed frequency-bias contours) at 12-h lead time, comparing
  the six-hourly MLP (blue) against the independence-assumption ensemble control (red).
  Panels show curves for six-hourly accumulated precipitation exceeding (a) 0.25\,mm,
  (b) 2.5\,mm, and (c) 10.0\,mm, with each curve built by cumulating contingency-table
  counts across forecast-probability bins from 0 to 100\% in 10\% increments.}
\label{fig:perfdiag_12h}
\end{figure}

\begin{figure}[H]
\centering
\includegraphics[width=\textwidth]{PerfDiagram_6h_MLP_MRMS_3panel_2025011000_to_2025123112_48h.png}
\caption{As in Fig.~\ref{fig:perfdiag_12h}, but for 48-h lead time.}
\label{fig:perfdiag_48h}
\end{figure}

Figure~\ref{fig:4panel_24h} shows the spatial structure of the six-hourly MLP
probabilistic forecasts for a representative Pacific Northwest US case.  The initial
condition is 12 UTC 4~December~2025 and the lead time is 36~h.  The GRAF
deterministic precipitation field (panel a) shows a broad precipitation region
over the Pacific Northwest, and the MLP probability fields (panels b--d) exhibit
terrain-related spatial structure that closely tracks the deterministic 
precipitation pattern.  The dot
marks the diagnostic location ($44^{\circ}$N, $122.25^{\circ}$W) used in
Fig.~\ref{fig:sample_36h}.

\begin{figure}[H]
\centering
\includegraphics[width=\textwidth]{Figs_6h/6h_MLP_4panel_IC2025120412_lead36h.png}
\caption{Four-panel plot of six-hourly MLP probabilistic forecasts for the initial
  condition of 12~UTC 4~December~2025 at 36-h lead time.  (a) GRAF deterministic
  six-hourly precipitation (mm); (b) probability of six-hourly accumulation
  $\geq$0.25\,mm; (c) probability $\geq$2.5\,mm; (d) probability $\geq$10\,mm.
  The dot marks the verification location ($44^{\circ}$N, $122.25^{\circ}$W)
  examined in Fig.~\ref{fig:sample_36h}.}
\label{fig:4panel_24h}
\end{figure}

Figure~\ref{fig:sample_36h} shows the predicted cumulative distribution function (CDF)
and exceedance probability at the marked location for the same initial condition
at 36-h lead time.  The six thin colored lines show the six consecutive hourly
distributions entering the MLP (lead hours 31--36), and the thick black line shows
the MLP six-hourly output distribution.  The MLP combines the information carried
across the six individual hours into a single coherent six-hourly predictive
distribution.  The six-hourly CDF is shifted well to the right of any single
hourly CDF, reflecting a sustained rain event in which appreciable
precipitation accumulates in most or all of the six hours rather than being
concentrated in one; summed over six hours, even a modest per-hour rate
compounds into a six-hourly total well beyond what any individual hour's
distribution would suggest.

\begin{figure}[H]
\centering
\includegraphics[width=\textwidth]{Figs_6h/6h_sample_IC2025120412_lead36h.png}
\caption{Predicted CDF (left panel) and exceedance probability (right panel) at
  $44^{\circ}$N, $122.25^{\circ}$W for the initial condition of 12~UTC
  4~December~2025 at 36-h lead time.  Thin colored lines show the six consecutive
  hourly gamma-mixture distributions (lead hours 31--36); the thick black line shows
  the MLP six-hourly output distribution.}
\label{fig:sample_36h}
\end{figure}

% =====================================================================
\section{Summary}
\label{sec:summary}
% =====================================================================

We have described a method for generating six-hourly probabilistic precipitation
forecasts from a sequence of six consecutive hourly probabilistic forecasts
produced by the Attention Residual U-Net.  The core of the approach is a
compact MLP with architecture $38 \to 72 \to 144 \to 72 \to 36 \to 12 \to 6$
that maps the six distributional parameters from each of six consecutive
hourly lead times to the six parameters of a six-hourly zero-inflated
two-component Gamma mixture.  Several design choices proved important to
this approach.  Operating on distributional parameters rather than raw
fields avoids storing full-resolution spatial feature data for training and
allows the six-hourly model to leverage the calibration properties already
embedded in the hourly forecasts.  Training uses importance sampling biased
toward grid points with higher probability of precipitation, combined with
a stricter MRMS quality threshold ($Q > 0.6$) than used for hourly training,
so that the training set emphasizes meteorologically relevant situations
and limits contamination from poor-quality radar retrievals.  The
output-layer is initialized climatologically via method-of-moments Gamma
fitting, which ensures stable convergence from the first epoch without
requiring EM-based two-component fitting.  Input features are Z-score
normalized, appropriate for the wide and \textit{a priori} unknown dynamic
range of the gamma-mixture parameters across lead times and seasons.
Finally, the model is trained with an NLL loss using log-sum-exp
stabilization and per-sample capping, providing a strictly proper scoring
rule that directly optimizes calibration of the predicted six-hourly
distribution.

Reliability diagrams for 12-h and 48-h lead times demonstrate that the
six-hourly MLP produces well-calibrated probabilistic forecasts for
precipitation thresholds of 0.25, 2.5, and 10.0\,mm, and is consistently
more reliable than a naive independence-assumption ensemble control that
sums independent Monte Carlo draws from the same hourly gamma-mixture
forecasts.  Performance diagrams built from the same contingency tables tell
a different, complementary story: the MLP's and the control's curves lie
close together at both lead times and across all three thresholds, with
maximum critical success index differing by 0.01 or less between the two
methods and the control's curve occasionally sitting marginally ahead of the
MLP's (most noticeably at the 10-mm threshold, 12-h lead time).  Together,
these results indicate that the MLP's demonstrated advantage over the naive
independence assumption lies specifically in probabilistic calibration --
the numeric trustworthiness of its forecast probabilities -- rather than in
its ability to discriminate events from non-events, since the two methods
consume the same hourly marginal information and therefore rank cases very
similarly.  An example case study plot showed a realistic-looking
relationship between precipitation amount and terrain relief.

Several aspects of the design were fixed by initial choice rather than by
systematic comparison, and are flagged here explicitly rather than left
implicit.  The two cyclic day-of-year features are motivated by the
expectation that a model trained on all twelve months, rather than a
handful of them, needs an explicit seasonal signal to distinguish, for
example, a summer convective regime from a winter frontal one; no ablation
comparing the MLP with and against these two features was performed, so
this expectation is plausible but unconfirmed, not demonstrated.  Likewise,
the specific hidden-layer widths $[72, 144, 72, 36, 12]$, the
expand-then-compress bottleneck shape, and the climatological method-of-
moments output-layer initialization (Sec.~\ref{sec:training}) were each
chosen once and used throughout, without a controlled comparison against
alternative widths, a monotonically shrinking (no bottleneck) architecture,
random initialization, or the EM-based two-component initialization used by
the hourly AttnResUNet.  It is therefore unknown how much of the reported
skill depends on these particular choices versus being robust to
reasonable variations in them.  The MLP architecture was not subject
to extensive hyperparameter search; a more thorough investigation of
hidden-layer widths or the use of residual connections within the MLP
might improve performance slightly.  Finally, incorporating additional predictors
beyond the hourly gamma-mixture parameters---such as GFS large-scale
moisture fields or time-of-day indicators---could potentially improve
the six-hourly distributional estimates, particularly for thresholds
associated with convective precipitation.

A method like this does not provide ensemble scenarios such as are
possible from ensembles of forecasts.  Such scenarios may be used in
ensemble streamflow predictions \citep{schaake2007hepex}.  This is an
admitted limitation of the method.  \citet{bellier2023} address a related
problem -- deriving coherent multi-temporal-scale precipitation guidance
from finer-resolution ensemble forecasts -- using a modulation mechanism
applied on top of standard univariate calibration and reordering methods
such as the Schaake shuffle; unlike the MLP described here, that approach
preserves ensemble scenarios by construction, at the cost of depending on
an underlying ensemble forecast rather than a single deterministic-model
probabilistic output.

A second admitted limitation concerns missing or degraded hourly input.
The MLP requires all six hourly gamma-mixture parameter sets spanning the
accumulation window to be present; every script in this pipeline that
assembles the six-hour input vector, whether for verification (as used
here) or for the standalone case-study plots, checks for the existence of
all six hourly files and simply declines to produce a six-hourly forecast
-- returning no result for that initialization and lead time, rather than
substituting a default, an interpolated, or a climatological value -- if
even one of the six is missing.  This all-or-nothing behavior is adequate
for retrospective verification, where incomplete cases can simply be
excluded, as was done here, but it is not yet a complete solution for
operational use.  A single missing or flagged-bad hourly input, for
example from a GRAF or GFS data outage upstream of the hourly AttnResUNet,
would presently cause the entire six-hourly forecast for that
initialization and lead time to be withheld.  No fallback strategy, such as
substituting a persisted forecast from an earlier hour, an hourly
climatology, or the mean of the remaining five hours, has been implemented
or evaluated, and this would need to be addressed before the method could
be relied upon in an operational setting where upstream data gaps do
occur.

% =====================================================================
% REFERENCES
% =====================================================================

\bibliographystyle{ametsoc2014}

\begin{thebibliography}{99}

\bibitem[Baran and Nem\'{o}da(2016)]{baran2016}
Baran, S., and D.\ Nem\'{o}da, 2016:
Censored and shifted Gamma distribution based EMOS model for probabilistic
quantitative precipitation forecasting.
\textit{Environmetrics}, \textbf{27}, 280--292,
doi:10.1002/env.2391.

\bibitem[Bellier \textit{et al.}(2023)]{bellier2023}
Bellier, J., B.\ Whitin, M.\ Scheuerer, J.~D.\ Brown, and T.~M.\ Hamill, 2023:
A multi-temporal-scale modulation mechanism for the postprocessing of
precipitation ensemble forecasts: Benefits for streamflow forecasting.
\textit{J.\ Hydrometeor.}, \textbf{24}, 659--673,
doi:10.1175/JHM-D-22-0119.1.

\bibitem[Ben Bouall\`{e}gue(2013)]{bouballegue2013}
Ben Bouall\`{e}gue, Z., 2013:
Calibrated short-range ensemble precipitation forecasts using extended
logistic regression with interaction terms.
\textit{Wea.\ Forecasting}, \textbf{28}, 515--524,
doi:10.1175/WAF-D-12-00062.1.

\bibitem[Bentzien and Friederichs(2012)]{bentzien2012}
Bentzien, S., and P.\ Friederichs, 2012:
Generating and calibrating probabilistic quantitative precipitation forecasts
from the high-resolution NWP model COSMO-DE.
\textit{Wea.\ Forecasting}, \textbf{27}, 988--1002,
doi:10.1175/WAF-D-11-00101.1.

\bibitem[Gasc\'{o}n \textit{et al.}(2019)]{gascon2019}
Gasc\'{o}n, E., D.\ Lavers, T.~M.\ Hamill, D.~S.\ Richardson,
Z.\ Ben Bouall\`{e}gue, M.\ Leutbecher, and F.\ Pappenberger, 2019:
Statistical postprocessing of dual-resolution ensemble precipitation forecasts
across Europe.
\textit{Q.\ J.\ R.\ Meteorol.\ Soc.}, \textbf{145}, 3218--3235,
doi:10.1002/qj.3615.

\bibitem[Ghazvinian \textit{et al.}(2020)]{ghazvinian2020}
Ghazvinian, M., Y.\ Zhang, and D.-J.\ Seo, 2020:
A nonhomogeneous regression-based statistical postprocessing scheme for
generating probabilistic quantitative precipitation forecasts.
\textit{J.\ Hydrometeor.}, \textbf{21}, 2275--2291,
doi:10.1175/JHM-D-20-0019.1.

\bibitem[Ghazvinian \textit{et al.}(2022)]{ghazvinian2022}
Ghazvinian, M., Y.\ Zhang, T.~M.\ Hamill, D.-J.\ Seo, and N.\ Fernando, 2022:
Improving probabilistic quantitative precipitation forecasts using short
training data through artificial neural networks.
\textit{J.\ Hydrometeor.}, \textbf{23}, 1365--1382,
doi:10.1175/JHM-D-22-0021.1.

\bibitem[Hamill(1999)]{hamill1999}
Hamill, T.~M., 1999:
Hypothesis tests for evaluating numerical precipitation forecasts.
\textit{Wea.\ Forecasting}, \textbf{14}, 155--167,
doi:10.1175/1520-0434(1999)014$<$0155:HTFENP$>$2.0.CO;2.

\bibitem[Hamill and Scheuerer(2018)]{hamill2018}
Hamill, T.~M., and M.\ Scheuerer, 2018:
Probabilistic precipitation forecast postprocessing using quantile mapping
and rank-weighted best-member dressing.
\textit{Mon.\ Wea.\ Rev.}, \textbf{146}, 4079--4098,
doi:10.1175/MWR-D-18-0147.1.

\bibitem[Hamill \textit{et al.}(2017)]{hamill2017nbm}
Hamill, T.~M., E.\ Engle, D.\ Myrick, M.\ Peroutka, C.\ Finan,
and M.\ Scheuerer, 2017:
The U.S.\ National Blend of Models for statistical postprocessing of
probability of precipitation and deterministic precipitation amount.
\textit{Mon.\ Wea.\ Rev.}, \textbf{145}, 3441--3463,
doi:10.1175/MWR-D-16-0331.1.

\bibitem[Hamill \textit{et al.}(2023)]{hamill2023nbm}
Hamill, T.~M., D.~R.\ Jesberg, and L.\ Smith, 2023:
Improving National Blend of Models probabilistic precipitation forecasts using
long time series of reforecasts and precipitation reanalyses.\ Part~I: Methods.
\textit{Mon.\ Wea.\ Rev.}, \textbf{151}, 1521--1534,
doi:10.1175/MWR-D-22-0308.1.

\bibitem[Hamill(2026)]{hamill2026a}
Hamill, T.~M., 2026:
High-resolution calibrated probabilistic hourly precipitation from a
deterministic forecast.
\textit{TBD}, \textbf{TBD}, TBD--TBD,
doi:TBD.

\bibitem[Hughes and Sangster(1979)]{hughes1979}
Hughes, L.~A., and W.~E.\ Sangster, 1979:
Combining precipitation probabilities.
\textit{Mon.\ Wea.\ Rev.}, \textbf{107}, 520--524,
doi:10.1175/1520-0493(1979)107$<$0520:CPP$>$2.0.CO;2.

\bibitem[Hu \textit{et al.}(2023)]{hu2023}
Hu, W., M.\ Ghazvinian, W.~E.\ Chapman, A.\ Sengupta, F.~M.\ Ralph, and
L.\ Delle Monache, 2023:
Deep learning forecast uncertainty for precipitation over the western
United States.
\textit{Mon.\ Wea.\ Rev.}, \textbf{151}, 1367--1385,
doi:10.1175/MWR-D-22-0268.1.

\bibitem[Ioffe and Szegedy(2015)]{ioffe2015}
Ioffe, S., and C.\ Szegedy, 2015:
Batch normalization: Accelerating deep network training by
reducing internal covariate shift.
\textit{Proc.\ 32nd Int.\ Conf.\ Mach.\ Learn.}, 448--456.

\bibitem[Jesberg \textit{et al.}(2023)]{jesberg2023}
Jesberg, D.~R., T.~M.\ Hamill, and L.\ Smith, 2023:
Improving National Blend of Models probabilistic precipitation forecasts using
long time series of reforecasts and precipitation reanalyses.\ Part~II: Results.
\textit{Mon.\ Wea.\ Rev.}, \textbf{151}, 1535--1550,
doi:10.1175/MWR-D-22-0310.1.

\bibitem[Kingma and Ba(2015)]{kingma2015}
Kingma, D.~P., and J.\ Ba, 2015:
Adam: A method for stochastic optimization.
\textit{Proc.\ 3rd Int.\ Conf.\ Learn.\ Represent.}

\bibitem[Messner \textit{et al.}(2014)]{messner2014}
Messner, J.~W., G.~J.\ Mayr, A.\ Zeileis, and D.~S.\ Wilks, 2014:
Heteroscedastic extended logistic regression for postprocessing of ensemble
guidance.
\textit{Mon.\ Wea.\ Rev.}, \textbf{142}, 448--456,
doi:10.1175/MWR-D-13-00271.1.

\bibitem[Nair and Hinton(2010)]{nair2010}
Nair, V., and G.~E.\ Hinton, 2010:
Rectified linear units improve restricted Boltzmann machines.
\textit{Proc.\ 27th Int.\ Conf.\ Mach.\ Learn.}, 807--814.

\bibitem[Roebber(2009)]{roebber2009}
Roebber, P.~J., 2009:
Visualizing multiple measures of forecast quality.
\textit{Wea.\ Forecasting}, \textbf{24}, 601--608,
doi:10.1175/2008WAF2222159.1.

\bibitem[Rumelhart \textit{et al.}(1986)]{rumelhart1986}
Rumelhart, D.~E., G.~E.\ Hinton, and R.~J.\ Williams, 1986:
Learning representations by back-propagating errors.
\textit{Nature}, \textbf{323}, 533--536,
doi:10.1038/323533a0.

\bibitem[Schaake \textit{et al.}(2007)]{schaake2007hepex}
Schaake, J.~C., T.~M.\ Hamill, R.\ Buizza, and M.\ Clark, 2007:
HEPEX: The Hydrological Ensemble Prediction Experiment.
\textit{Bull.\ Amer.\ Meteor.\ Soc.}, \textbf{88}, 1541--1547,
doi:10.1175/BAMS-88-10-1541.

\bibitem[Scheuerer(2014)]{scheuerer2014}
Scheuerer, M., 2014:
Probabilistic quantitative precipitation forecasting using ensemble model
output statistics.
\textit{Q.\ J.\ R.\ Meteorol.\ Soc.}, \textbf{140}, 1086--1096,
doi:10.1002/qj.2183.

\bibitem[Scheuerer and Hamill(2015)]{scheuerer2015}
Scheuerer, M., and T.~M.\ Hamill, 2015:
Statistical postprocessing of ensemble precipitation forecasts by fitting
censored, shifted Gamma distributions.
\textit{Mon.\ Wea.\ Rev.}, \textbf{143}, 4578--4596,
doi:10.1175/MWR-D-15-0061.1.

\bibitem[Skamarock \textit{et al.}(2012)]{skamarock2012}
Skamarock, W.~C., J.~B.\ Klemp, M.~G.\ Duda, L.~D.\ Fowler,
S.\ Park, and T.~D.\ Ringler, 2012:
A multiscale nonhydrostatic atmospheric model using centroidal
Voronoi tesselations and C-grid staggering.
\textit{Mon.\ Wea.\ Rev.}, \textbf{140}, 3090--3105.

\bibitem[Sloughter \textit{et al.}(2007)]{sloughter2007}
Sloughter, J.~M., A.~E.\ Raftery, T.\ Gneiting, and C.\ Fraley, 2007:
Probabilistic quantitative precipitation forecasting using Bayesian model
averaging.
\textit{Mon.\ Wea.\ Rev.}, \textbf{135}, 3209--3220,
doi:10.1175/MWR3441.1.

\bibitem[Voisin \textit{et al.}(2010)]{voisin2010}
Voisin, N., J.~C.\ Schaake, and D.~P.\ Lettenmaier, 2010:
Calibration and downscaling methods for quantitative ensemble precipitation
forecasts.
\textit{Wea.\ Forecasting}, \textbf{25}, 1603--1627,
doi:10.1175/2010WAF2222367.1.

\bibitem[Wilks(1990)]{wilks1990}
Wilks, D.~S., 1990:
On the combination of forecast probabilities for consecutive precipitation
periods.
\textit{Wea.\ Forecasting}, \textbf{5}, 640--650,
doi:10.1175/1520-0434(1990)005$<$0640:OTCOFP$>$2.0.CO;2.

\bibitem[Wilks(1991)]{wilks1991}
Wilks, D.~S., 1991:
Representing serial correlation of meteorological events and forecasts in
dynamic decision-analytic models.
\textit{Mon.\ Wea.\ Rev.}, \textbf{119}, 1640--1662,
doi:10.1175/1520-0493(1991)119$<$1640:RSCOME$>$2.0.CO;2.

\bibitem[Zhang and Coauthors(2016)]{zhang2016}
Zhang, J., and Coauthors, 2016:
Multi-Radar Multi-Sensor (MRMS) quantitative precipitation
estimation: Initial operating capabilities.
\textit{Bull.\ Amer.\ Meteor.\ Soc.}, \textbf{97}, 621--638.

\bibitem[Zhang \textit{et al.}(2017)]{zhang2017}
Zhang, J., H.\ Wu, M.\ Scheuerer, J.\ Schaake, and C.\ Kongoli, 2017:
Comparison of probabilistic quantitative precipitation forecasts from two
postprocessing mechanisms.
\textit{J.\ Hydrometeor.}, \textbf{18}, 2873--2891,
doi:10.1175/JHM-D-16-0293.1.

\end{thebibliography}

\end{document}
